% GENERATED FILE. Edit canonical lessons, then run npm run generate:books.
% canonical-source-hash: fnv1a64:07135320fef657d6
% canonical-lessons: TE-C152-jhari, TE-S173-letter-jha, TE-R152-jhari-recall

\chapter{A Word for a Stream — Then Its First Letter}
\label{ch:te-jhari-jha}
\hlchaptermodality{152}

\begin{chapteropening}
\textbf{By the end of this chapter:} \emph{I can say a Telugu word for a stream, recognise its uncommon first letter, and write that letter in a source-backed order.}

Take \te{ఝ} from the already-known word \te{ఝరి} and write its five movements in five pen-down runs.
\end{chapteropening}

\section[\texorpdfstring{\te{ఝరి} (jhari)}{ఝరి (jhari)}]{\te{ఝరి} (jhari) — a stream}

\label{lesson:TE-C152-jhari}



\begin{quote}
What did \textbf{\te{వాఙ్మయం}} name? (\textquotedblleft{}Literature.\textquotedblright{}) \emph{Find:} the familiar \textbf{\te{రి}} at the end of today's word — \textbf{\te{ర}} carries the short \textbf{\te{ి}} sign

Picture a narrow stream moving down a hillside. This lesson gives that stream a name before it asks you to study the uncommon first letter.
\end{quote}

\subsection*{You'll want to know: \te{ఝరి}}
\textbf{\te{ఝరి}} — \emph{jhari} — \textquotedblleft{}a stream\textquotedblright{} or \textquotedblleft{}mountain stream\textquotedblright{}.

Say it in two gentle pieces: \emph{jha · ri}. Say the word as \textbf{\te{ఝ}-\te{రి}}. The first sound is a breathy \textbf{jh}; the second piece is the familiar \textbf{\te{రి}}.

\textbf{\te{ఝరి}} is an authentic learned or literary word shared with Sanskrit. It is not an everyday word you must force into conversation. Keep its concrete image today: a stream running down a slope.

\begin{grammarlens}[title={What you've built}]
Meaning comes first. You can recognise and say the whole word before the next lesson asks your hand to take its first shape apart.
\end{grammarlens}

\subsection*{Your turn}
\begin{itemize}
\raggedright
  \item \emph{Say it:} \emph{jha · ri}, slowly
  \item \emph{Say it:} \emph{jhari}, as one word
  \item \emph{Read it:} \textbf{\te{ఝరి}} — a stream
  \item \emph{Find:} the familiar \textbf{\te{రి}} at the end
\end{itemize}

\begin{tcolorbox}[breakable,colback=teal!4,colframe=teal!35!black,title={Before you move on}]
What does \textbf{\te{ఝరి}} mean? (\textquotedblleft{}A stream\textquotedblright{} or \textquotedblleft{}mountain stream.\textquotedblright{}) Say it once more.
\end{tcolorbox}
\section[\texorpdfstring{\te{ఝ} (jha)}{ఝ (jha)}]{\te{ఝ} — the breathy first letter in \te{ఝరి}}

\label{lesson:TE-S173-letter-jha}



\begin{quote}
One gentle look back first: \textbf{\te{వాఙ్మయం}} means \textquotedblleft{}literature,\textquotedblright{} and its uncommon letter is \textbf{\te{ఙ}}. Point to \textbf{\te{ఙ}} without copying it. Today you will keep that shape distinct from a different uncommon letter.

Now say \textbf{\te{ఝరి}} once. What does it mean? (\textquotedblleft{}A stream\textquotedblright{} or \textquotedblleft{}mountain stream.\textquotedblright{})

Now slow the first piece: \emph{jha}. That breathy \textbf{jh} base letter is today's whole task.
\end{quote}

\subsection*{Script you'll notice: \te{ఝ}}
\textbf{\te{ఝ}} — \emph{jha}.

In \textbf{\te{ఝరి}}, the first base letter is \textbf{\te{ఝ}}. The rest, \textbf{\te{రి}}, is already familiar. Read the whole word in two pieces: \textbf{\te{ఝ}-\te{రి}}.

This aspirated letter is uncommon in everyday Telugu vocabulary. That is why the book gave you the authentic word \textbf{\te{ఝరి}} before asking you to practise the shape.

\subsection*{Writing: \te{ఝ} — five movements, five runs}
\hlblockfigure{\detokenize{figures/TE-S173-letter-jha-filmstrip.pdf}}{How \te{ఝ} is written, stroke by stroke}

Follow the filmstrip slowly. Movement 1 circles the broad left bowl. Lift. Movement 2 circles the middle bowl. Lift again. Movement 3 circles the right bowl. Movement 4 sweeps through the separate upper flourish. Movement 5 draws the separate stem downward below the middle.

The numbered route is one attested school-style order; Telugu handwriting can vary. Keep all five pen-down runs distinct and let the three bowls stay round.

\subsection*{Your turn}
\begin{itemize}
\raggedright
  \item \emph{Trace it:} \textbf{\te{ఝ}} once, following the five numbered movements
  \item \emph{Copy:} \textbf{\te{ఝ}} twice without tracing
  \item \emph{Build:} \textbf{\te{ఝ}} + \textbf{\te{రి}} = \textbf{\te{ఝరి}}
  \item \emph{Read it:} \textbf{\te{ఝరి}}
\end{itemize}

\begin{tcolorbox}[breakable,colback=teal!4,colframe=teal!35!black,title={Before you move on}]
Which letter is this — \textbf{\te{ఝ}}? What word gave it to you? (\textbf{\te{ఝరి}}.)
\end{tcolorbox}
\section[\texorpdfstring{\te{ఝరి} (jhari)}{ఝరి (jhari)}]{\te{ఝరి} — meaning and shape together}

\label{lesson:TE-R152-jhari-recall}



\begin{quote}
Picture a stream running down a hillside. Say the Telugu word for it. (\textbf{\te{ఝరి}}.)
\end{quote}

\subsection*{Your turn}
\emph{Read it:} \textbf{\te{ఝరి}} once \emph{Cover:} it \emph{Write it:} the whole word from memory, with only the picture of a stream as your cue

\emph{Check:} the uncovered model — both the whole word and its first base letter \emph{Write it:} a circle round \textbf{\te{ఝ}} — then check its three rounded bowls, upper flourish, and separate downward stem, keeping the five pen-down runs distinct

\begin{itemize}
\raggedright
  \item \emph{Write it:} \textbf{\te{ఝరి}} once from the picture cue
  \item \emph{Circle:} the first base letter \textbf{\te{ఝ}}
  \item \emph{Say it:} “a stream”
\end{itemize}

\begin{tcolorbox}[breakable,colback=teal!4,colframe=teal!35!black,title={Before you move on}]
\emph{Write it:} without looking back, the Telugu word for “a stream,” then a circle round its first base letter

(\textbf{\te{ఝరి}}; the circle goes round \textbf{\te{ఝ}}.)
\end{tcolorbox}
